% ─────────────────────────────────────────────────────────────
\phantomsection
\chapter{Conclusion}
\label{sec:conclusion}

This thesis investigated whether persistent homology, applied to short wearable
accelerometer windows via a delay-embedding pipeline, can support recall-oriented
fall detection that is simultaneously rigorous in its evaluation, principled in its
feature construction, and lightweight enough for edge deployment.  The investigation
was organised around four research hypotheses and carried out on three public
benchmark datasets---MobiFall, SisFall, and FAD\_40Hz---under a three-protocol
subject-aware evaluation framework.  The findings are summarised below at three
levels of description: mathematical, statistical, and engineering.

% ─────────────────────────────────────────────────────────────
\section{Mathematical Perspective}
\label{subsec:conclusion-math}

\paragraph{Representation effectiveness (H1).}
The feature map $\phi_\lambda\colon\R^{n\times 3}\to\R^{24}$, with hyperparameters
$\lambda^*_K$ identified by Bayesian optimisation, achieves mean subject-general
$F_2$ of $0.9709$ (LR) and $0.9700$ (SVM) on SisFall, $0.9429$ (LR) and
$0.9592$ (SVM) on FAD\_40Hz, and $0.9049$ (LR) and $0.8918$ (SVM) on MobiFall
under the leave-one-subject-out protocol.  All values exceed the pre-specified
thresholds ($\geq 0.93$ on SisFall and FAD\_40Hz; $\geq 0.80$ on MobiFall) for at
least one classifier, and recall is above $0.95$ across all dataset-classifier
combinations.  H1 is conditionally confirmed: the representation meets its own
targets, subject to the caveat that no baseline comparison has been conducted.

The theoretical basis for this effectiveness is the Cohen--Steiner stability
theorem, which guarantees that the bottleneck distance between two persistence
diagrams is bounded by the supremum norm of the perturbation between the
corresponding input functions.  This bound is a mathematical theorem, not an
empirical observation, and it provides a formal noise-robustness certificate that
no competing feature family---handcrafted statistical features, spectral
coefficients, or deep-network embeddings---possesses in comparable form.

\paragraph{Dataset-specific topological geometry (H2).}
The three datasets converged to structurally distinct optimal configurations:
MobiFall selected a SparseRips complex at a 5.0~s window with Chebyshev distance,
SisFall selected an Alpha complex at a 1.0~s window, and FAD\_40Hz selected a
standard Rips complex at a 1.0~s window with Manhattan distance.  Every dimension of
the hyperparameter tuple $\lambda$ varies across at least two datasets.  H2 is
qualitatively confirmed: no universal configuration exists for these three datasets.

Mathematically, this finding is consistent with the geometry of the underlying
delay-embedded point clouds.  The Alpha complex is derived from the Euclidean
Delaunay triangulation and is appropriate when the point cloud is convex and
well-conditioned; SisFall's controlled, elderly-subject falls produce exactly this
geometry.  The SparseRips complex handles large, noisy clouds efficiently by
subsampling; MobiFall's continuous-wear sensor data requires this treatment.  The
differing window lengths encode a genuine difference in fall dynamics: MobiFall falls
require a 5.0~s window to capture the pre-impact, impact, and post-fall phases
jointly, while SisFall and FAD\_40Hz falls are discriminable in 1.0~s.

\paragraph{Representation dominates classifier choice (H3).}
The absolute difference in General-protocol $F_2$ between LR and SVM is $0.0131$
on MobiFall, $0.0009$ on SisFall, and $0.0163$ on FAD\_40Hz---all below $0.02$.
By contrast, the $F_2$ range across 200 Optuna trials is $0.179$ on MobiFall,
$0.053$ on SisFall, and $0.108$ on FAD\_40Hz.  The ratio of $\lambda$-tuning
variation to classifier gap is $13.7\times$ on MobiFall, $59\times$ on SisFall,
and $6.6\times$ on FAD\_40Hz.

H3 is strongly confirmed.  The mathematical interpretation is that $\phi_{\lambda^*}$
with the optimal representation transforms the nonlinear fall-versus-ADL boundary
into an approximately linearly separable problem in $\R^{24}$.  Once this
transformation is achieved, both logistic regression (linear boundary) and SVM
(maximum-margin boundary, with or without an RBF kernel) operate near their
theoretical limit, producing near-equivalent results.  The signal resides in the
map $\phi_{\lambda^*}$, not in the classifier class.

% ─────────────────────────────────────────────────────────────
\section{Statistical Perspective}
\label{subsec:conclusion-stat}

\paragraph{Threshold personalisation as a distributional phenomenon (H4).}
The subject-level sweep gain $\Delta F_2(s) = F_2(\tau^*(s),s) - F_2(\hat\tau_0(s),s)$,
where $\hat\tau_0(s)$ is the PR-curve trained threshold, has a right-skewed
distribution in all six dataset-classifier combinations.  The mean gains are
$0.035$--$0.047$ on MobiFall, $0.025$--$0.035$ on FAD\_40Hz, and
$0.004$--$0.008$ on SisFall.  The coefficient of variation exceeds $50\%$ in all
cases, reaching $138\%$ for SisFall LR---indicating that the distribution is highly
dispersed relative to its mean.

H4 is confirmed in a refined form.  The proposal predicted that the distribution
of $\Delta_s$ would contain both positive and negative values.  In practice, all
gains are non-negative by construction (the sweep is an argmax), but the
distribution is far from uniform: it concentrates near zero for the majority of
subjects while a small minority achieves gains exceeding $0.10$ units of $F_2$.
Subject~11 on FAD\_40Hz (LR) gains $+0.1165$, lifting $F_2$ from $0.8046$ to
$0.9211$ by adjusting the threshold alone.  This heterogeneity is incompatible with
a uniform-shift explanation: if the classifier's posterior probabilities were
globally miscalibrated in the same direction for all subjects, a single global
threshold offset would equalise all gains.  The observed distribution implies that
the optimal decision boundary location is genuinely subject-specific.

\paragraph{Evaluation integrity under LOSO.}
The leave-one-subject-out protocol with $R=15$ independent class-balancing
repetitions separates two sources of variability that are conflated in random
train-test splits: (i) within-subject repetition variance, captured by the
repetition-level $t$-confidence intervals, and (ii) across-subject heterogeneity,
visible in the per-subject tables.  The Naive protocol consistently produces results
comparable to or marginally lower than the General protocol on SisFall (Naive
$F_2=0.9713$ vs.\ General $0.9709$ for LR), suggesting that LOSO is not
materially penalised by train-test class mismatch on large, diverse datasets.  On
MobiFall, where the subject pool is small (9 subjects), the repetition-level
confidence intervals are wider ($\pm0.039$ for LR), correctly reflecting the
limited statistical power of a 9-fold LOSO evaluation.

The Personal protocol reveals a further finding: full per-subject model training
degrades $F_2$ severely on MobiFall ($\Delta = -0.172$ for LR) but only modestly on
SisFall ($\Delta = -0.030$ for LR), with the magnitude of degradation inversely
correlated with the subject pool size.  This confirms that personalised model
training is only beneficial when sufficient per-subject data is available, and that
the threshold sweep is the appropriate personalisation mechanism when data is sparse.

% ─────────────────────────────────────────────────────────────
\section{Engineering Perspective}
\label{subsec:conclusion-eng}

\paragraph{Inference-time personalisation without retraining.}
The central engineering contribution of the thesis is the identification of
decision-threshold optimisation as a first-class personalisation mechanism that
is qualitatively different from all prior alternatives.  The threshold
$\tau\in\mathcal{T}=\{0.05,\ldots,0.95\}$ is the only parameter that can be
adjusted at inference time, without gradient computation, without access to the
training data, and without transmission of data to a server.  The computational
cost of a full per-subject sweep is $O(|\mathcal{T}|\cdot|W_s|)$ comparisons,
amounting to less than one second of wall time for all subjects in any of the three
datasets.  The ratio to the Bayesian optimisation cost exceeds $10^3$.

This asymmetry has a direct consequence for product design.  A wearable fall
detection system based on the proposed pipeline can be shipped with a population-level
model ($\phi_{\lambda^*}$, $f_\theta$) and personalised to each user at installation
time by sweeping the threshold over a small held-out calibration set.  No server
round-trip is required.  No labeled falls are required from the target user in the
first instance---the sweep can be initialised from simulated events during setup and
refined as real events accumulate.  This is the only personalisation mechanism among
those surveyed in this thesis that simultaneously satisfies the zero-retraining,
privacy-preserving, and inference-time requirements.

\paragraph{CPU-only deployment profile.}
The entire pipeline---delay embedding, simplicial filtration, boundary matrix
reduction, feature vectorization, linear classification, and threshold sweep---is
executable on commodity CPU hardware.  Persistent homology computation on the
window sizes and embedding dimensions considered here ($m \leq 5$,
$|P_\lambda| \leq 150$ points) completes in tens of milliseconds per window on a
desktop-class processor core.  The Bayesian optimisation consumed ${\sim}75$~min
of CPU time on a 40-core HPC node; this cost is paid once per dataset, not per
user.

The ARM Cortex-A processors present in current mid-range smartphones are
comfortably capable of executing the per-window inference pipeline in real time at
50~Hz.  Modern smartwatch SoCs are more constrained, but the per-window
computation---feature extraction plus a matrix-vector multiply---is within their
budget.  The absence of a GPU requirement eliminates the most significant barrier
to embedded deployment and makes the pipeline uniquely well-suited to the
battery-constrained wearable context.

\paragraph{Failure mode analysis and minimum calibration set size.}
The most significant failure mode identified in this thesis is the collapse of
the Personal protocol on MobiFall: $F_2$ drops from $0.9049$ (General) to
$0.7325$ (Personal) for LR, and from $0.8918$ to $0.5981$ for SVM.  The cause is
statistical rather than algorithmic: MobiFall contains 9 subjects with
approximately 20--30 fall instances per subject, which is insufficient for stable
within-subject classifier estimation.  The implication for system designers is
concrete: a minimum calibration set size is required before full personal model
training is beneficial.  For the datasets studied here, that threshold lies somewhere
between the 9 subjects of MobiFall and the 13 subjects of FAD\_40Hz.  Below this
threshold, the threshold sweep is the safer personalisation mechanism.

This failure mode is not a weakness of TDA specifically; it would affect any
classifier trained on sparse per-subject data.  It is, however, an important
calibration of expectations for real deployment: the naive assumption that
``more personalization is always better'' is empirically refuted by the MobiFall
Personal protocol results.

\paragraph{Path to consumer deployment.}
Taken together, the three engineering findings---inference-time personalisation,
CPU-only operation, and the identified minimum calibration requirement---define a
concrete deployment pathway:

\begin{enumerate}
  \item A population-level model ($\phi_{\lambda^*}$, $f_\theta$) is trained once
    on a large labeled dataset and shipped with the device.
  \item At installation, the user performs a brief simulated fall session
    (or the device waits to accumulate confirmed fall events from normal wear).
    The per-subject threshold $\tau^*(s)$ is identified from the calibration
    scores in under one second.
  \item Subsequent detections use $\tau^*(s)$.  As more events accumulate, the
    threshold estimate can be refined.
  \item No server communication is required for any of steps 1--3; all computation
    is on-device.
\end{enumerate}

The natural next step beyond threshold personalisation is one-class anomaly detection:
training the model exclusively on the user's own activities of daily living,
treating falls as out-of-distribution events without requiring any labeled fall data.
The Cohen--Steiner stability of the TDA features ensures that the ADL manifold in
$\R^{24}$ is robustly estimated from normal wear data, and the threshold sweep
mechanism extends directly to the standoff distance from the one-class boundary.
This approach eliminates the cold-start problem entirely and represents the most
direct path from the current thesis to a consumer-ready, zero-label deployment.

% ─────────────────────────────────────────────────────────────
\section{Limitations and Open Questions}
\label{subsec:conclusion-limits}

Four limitations of the current work constrain the strength of the conclusions.

First, the absence of a baseline comparison means that H1 is validated only against
its own pre-specified targets.  The claim that TDA is superior to statistical-feature
or threshold-based baselines on these datasets cannot be made from the current
evidence, and establishing this comparison is the highest-priority open task.

Second, the cross-application drop matrix $M_{K,K'}$ was not computed.  H2 is
therefore confirmed only qualitatively (distinct optimal configurations) and not
quantitatively (measured performance penalty from transplanting $\lambda^*_K$ to
dataset $K'$).

Third, no ablation study was conducted.  The 24-dimensional feature vector
$\phi_\lambda$ contains $H_0$, $H_1$, and signal statistics; the relative
contribution of each subvector is unknown, and it is possible that a smaller
sub-map would achieve comparable performance with lower computational cost.

Fourth, all experiments assume a fixed sensor placement at the waist or in the
trouser pocket.  The stability theorem guarantees noise robustness within a
placement, but does not extend to fundamentally different placements (wrist, chest,
thigh).  Robustness to placement variation is an important practical question that
remains unaddressed.

% ─────────────────────────────────────────────────────────────
\section{Closing Remarks}

This thesis set out to answer a focused question: can a TDA pipeline for wearable
fall detection be simultaneously principled, lightweight, and subject-aware?  The
experimental evidence across three datasets and two classifiers supports an
affirmative answer on all three dimensions.  The pipeline is principled in the sense
that its key property---feature stability---is guaranteed by a mathematical theorem.
It is lightweight in the sense that it runs without GPU acceleration at latencies
compatible with real-time embedded operation.  It is subject-aware in the sense that
the decision threshold can be personalised at inference time, at negligible
computational cost, without retraining.

The broader lesson from H3 is perhaps the most transferable finding.  The gap between
a well-tuned and a poorly tuned topological representation ($0.10$--$0.18$ units of
$F_2$) dwarfs the gap between the best and worst classifier at the optimal
representation ($< 0.02$ units).  In the fall detection context, and likely in other
wearable time-series classification problems where the discriminative signal is
encoded in short-window shape geometry, choosing the right representation is the
dominant engineering decision.  The choice of classifier is, in an important sense,
secondary.

